\section*{Глава 1. Анализ предметной области}
\addcontentsline{toc}{section}{Глава 1. Анализ предметной области}
\stepcounter{section}

\subsection{Обзор существующих систем электронного документооборота}

На современном этапе цифровизации организаций системы электронного документооборота (СЭД) являются ключевым инструментом автоматизации управленческой деятельности. Они обеспечивают централизованное хранение документов, контроль версий, маршрутизацию согласования и разграничение прав доступа. В таблице~\ref{tab:competitors} представлен сравнительный анализ наиболее распространённых решений.

\begin{table}[ht!]
\centering
\caption{Сравнительный анализ систем электронного документооборота}
\label{tab:competitors}
\setlength{\tabcolsep}{3pt}
\small
\renewcommand{\arraystretch}{1.15}
\begin{tabularx}{\textwidth}{|p{3cm}|*{3}{>{\raggedright\arraybackslash}X|}}
\hline
\textbf{Критерий сравнения} & \textbf{1С:Документооборот} & \textbf{Directum} & \textbf{Alfresco} \\ \hline
\textbf{Платформа и доступ} & Desktop/Web, интеграция с 1С & Web, SaaS и On-premise & Web, кроссплатформенная \\ \hline
\textbf{Управление документами} & Полный цикл работы с документами, шаблоны, маршруты & Развитая маршрутизация и BPM & Гибкое хранение и версионирование \\ \hline
\textbf{Ролевая модель} & Гибкая настройка ролей и прав & Разграничение доступа и бизнес-роли & Настраиваемая модель доступа \\ \hline
\textbf{Аналитика и отчёты} & Расширенные отчёты и интеграция с учетными системами & BPM-аналитика процессов & Ограниченные встроенные средства \\ \hline
\textbf{Гибкость и кастомизация} & Требует доработок под специфику & Высокая, но сложная настройка & Высокая (open-source) \\ \hline
\textbf{Ключевой недостаток} & Высокая стоимость внедрения и сопровождения & Сложность внедрения для малых организаций & Требует технической доработки и поддержки \\ \hline
\end{tabularx}
\end{table}

Анализ показал, что существующие системы ориентированы преимущественно на крупные организации и корпоративный сектор. Они обладают высокой функциональной сложностью и требуют значительных ресурсов на внедрение, сопровождение и обучение сотрудников работе с ними. Для культурно-досуговых учреждений, таких как МУ ЦКиД «Факел» г. Фрязино, характерны более простые, но специфические процессы документооборота, а также небольшие бюджеты на поддержание этих процессов. Это обосновывает необходимость разработки легковесного и адаптированного веб-приложения, ориентированного на реальные потребности пользователей.

\subsection{Анализ программных средств и технологий разработки веб-приложений}

Выбор технологического стека определялся требованиями к простоте разработки, скорости внедрения, масштабируемости и удобству сопровождения. Архитектура приложения строится по клиент-серверной модели. В таблицах~\ref{tab:frontend},~\ref{tab:backend},~\ref{tab:database} представлен сравнительный анализ различных технологий в рамках создания веб-приложения для автоматизации процессов документооборота в культурно-досуговых учреждениях.

\begin{table}[ht!]
\centering
\caption{Сравнительный анализ клиентских технологий}
\label{tab:frontend}
\setlength{\tabcolsep}{3pt}
\small
\renewcommand{\arraystretch}{1.15}
\begin{tabularx}{\textwidth}{|p{3cm}|*{3}{>{\raggedright\arraybackslash}X|}}
\hline
\textbf{Технология} & \textbf{Преимущества} & \textbf{Недостатки} & \textbf{Итог} \\ \hline
\textbf{Чистый JS + HTML + CSS} & Минимальный порог входа, простота разработки, отсутствие сложной сборки & Ограниченные возможности масштабирования интерфейса & \textbf{Выбран} - оптимален для учебного проекта и умеренной сложности UI \\ \hline
\textbf{React} & Компонентная архитектура, масштабируемость, развитая экосистема & Усложнение проекта, необходимость настройки сборки & Не выбран - избыточен для текущих требований \\ \hline
\textbf{Vue.js} & Простота интеграции, реактивность, удобство разработки & Меньшая экосистема по сравнению с React & Не выбран - не требуется при простой логике интерфейса \\ \hline
\end{tabularx}
\end{table}

\begin{table}[ht!]
\centering
\caption{Сравнительный анализ серверных технологий}
\label{tab:backend}
\setlength{\tabcolsep}{3pt}
\small
\renewcommand{\arraystretch}{1.15}
\begin{tabularx}{\textwidth}{|p{3cm}|*{3}{>{\raggedright\arraybackslash}X|}}
\hline
\textbf{Технология} & \textbf{Преимущества} & \textbf{Недостатки} & \textbf{Итог} \\ \hline
\textbf{FastAPI} & Высокая производительность, асинхронность, автоматическая документация API, простота разработки & Требует разделения на API и frontend & \textbf{Выбран} - оптимален по балансу простоты и производительности \\ \hline
\textbf{Django} & «Из коробки» включает ORM, админ-панель, авторизацию & Избыточность и монолитность для небольшой системы & Не выбран - избыточен для проекта \\ \hline
\textbf{Node.js + Express} & Единый язык с frontend, гибкость & Требует ручной организации архитектуры, нет автоматической документации & Не выбран - нет критического преимущества перед FastAPI \\ \hline
\end{tabularx}
\end{table}
\newpage
\begin{table}[ht!]
\centering
\caption{Сравнительный анализ систем управления базами данных}
\label{tab:database}
\setlength{\tabcolsep}{3pt}
\small
\renewcommand{\arraystretch}{1.15}
\begin{tabularx}{\textwidth}{|p{3cm}|*{3}{>{\raggedright\arraybackslash}X|}}
\hline
\textbf{СУБД} & \textbf{Преимущества} & \textbf{Недостатки} & \textbf{Итог} \\ \hline
\textbf{PostgreSQL} & Надёжность, транзакции, сложные запросы, поддержка связей & Требует проектирования схемы & \textbf{Выбран} - оптимален для структурированных данных \\ \hline
\textbf{MySQL} & Простота и распространённость & Ограниченная функциональность сложных запросов & Не выбран - уступает PostgreSQL \\ \hline
\textbf{MongoDB} & Гибкая схема, масштабируемость & Сложность обеспечения связей и целостности & Не выбран - не подходит для строгой структуры данных \\ \hline
\end{tabularx}
\end{table}

Для реализации клиентской части приложения используются базовые веб-технологии: HTML, CSS и JavaScript. Данный подход позволяет реализовать пользовательский интерфейс без использования тяжеловесных фреймворков, обеспечивая простоту разработки и достаточную функциональность для решения поставленных задач. В процессе разработки использовались официальные спецификации и руководства по веб-стандартам~\cite{mdn_web_docs}.

Для реализации серверной части приложения выбрана платформа FastAPI, являющаяся современным фреймворком для разработки веб-приложений на языке Python. FastAPI обеспечивает высокую производительность за счёт асинхронной обработки запросов, а также предоставляет встроенную валидацию данных и автоматическую генерацию документации API~\cite{fastapi_docs}. Использование данного фреймворка позволяет эффективно реализовать REST-интерфейс для взаимодействия клиентской и серверной частей приложения.

Для хранения данных в приложении выбрана система управления базами данных PostgreSQL. Данная СУБД обеспечивает надежное хранение структурированных данных, поддержку транзакций и строгую целостность данных. Кроме того, PostgreSQL предоставляет широкие возможности для выполнения сложных запросов, индексации и масштабирования, что делает её подходящей для реализации систем документооборота~\cite{postgres_docs}. Использование реляционной модели данных позволяет эффективно организовать хранение сущностей приложения, таких как пользователи, документы, роли и права доступа, а также установить между ними необходимые связи.

\subsection{Формулировка цели и задач работы}

На основе проведенного анализа существующих решений (см. подраздел 1.1) и выбора технологического стека (см. подраздел 1.2) была сформулирована цель исследования.

\textbf{Целью работы} является проектирование и разработка веб-приложения для автоматизации документооборота в культурно-досуговых учреждениях на примере МУ ЦКиД «Факел» г. Фрязино.

Для достижения поставленной цели необходимо решить следующие \textbf{задачи}:
\begin{enumerate}
\item Провести анализ предметной области и выявить особенности документооборота в культурно-досуговых учреждениях на примере МУ ЦКиД «Факел».
\item Выполнить сравнительный анализ существующих систем электронного документооборота.
\item Спроектировать архитектуру клиент-серверного приложения и структуру базы данных.
\item Реализовать серверную часть приложения с использованием FastAPI и обеспечить обработку запросов пользователей.
\item Разработать клиентский интерфейс с использованием HTML, CSS и JavaScript.
\item Реализовать CRUD-операции для работы с документами и пользователями, а также систему разграничения прав доступа.
\item Провести тестирование разработанной системы и оценить её эффективность.
\end{enumerate}